
\section{What True Edification Is}

Installment 4: “Folkebladet,” March 15, 1916

Through the little I have said here about this matter, I have tried to indicate how indefinite and how lacking in a clear aim it is when prayer is offered for “salvation” for a soul.

And I believe that the same is true when prayer is offered for “edification for the believers.” For what is commonly meant by the expression “edification” is that the single person—the individual—may during the meeting be brought into a pleasant feeling of joy and satisfaction with himself or with his Christianity. In other words: joy and a sense of happiness in being a believer, in contrast to “the world.”

Now I do not mean to say that joy over salvation is unwarranted and harmful. On the contrary. I maintain, with God’s Word, that “joy in the Lord is our strength.” But in the first place, it is far from true that “joy in the Lord” is the same as a feeling of pleasant satisfaction with oneself and one’s conversion. No, “joy in the Lord” is the soul’s blessed rest—not in its repentance, conversion, faith, or sanctification, but in the completed salvation through Jesus Christ—procured for all, but received and availed of only by those who bow themselves to God with an undivided mind.

Now it is essentially true that the soul which “believes,” which receives in Jesus Christ salvation from sin and perdition, possesses something that “the world” does not. It is true that where surrender to God has taken place, a great change has occurred—the greatest that can happen to a human being. But this change, by which a person becomes willing to receive salvation, is a work of God himself through his Spirit and Word—a work which the Spirit has accomplished amid much resistance and many hindrances on our part.

“Joy in the Lord” cannot, therefore, be joy over the fact that “I have been converted,” but joy that salvation has been procured for me through Jesus’ atoning death; and when, through the Spirit’s guidance and working, I have learned to make it my own, to rest in it, and to build my whole blessedness upon it—then I possess something of the joy in the Lord. But the others—“the world”—can also receive a share in this “salvation.” For “God wills that all people should be saved and come to the knowledge of the truth.”

That the believers are edified is not the same as being confirmed in the thought that I am so thoroughly converted, or that in my heart I am certain that, if everyone were like me, all would be well in the world. No, “edification” is that, through the working of the Word, I am made to see how poor a Christian I am, how meager both faith and life are—in other words: that my knowledge of sin is deepened, my gratitude to God for salvation increased, and my zeal for the salvation of others made more ardent.

“Edification,” therefore, is this: that I receive more of the mind of Jesus, so that for me, as for him, it becomes “my food,” the wealth and happiness of my life, the secret fountain of joy in my soul, to be permitted to fight in self-denial, to take part in the Lord’s army in the battle against the evil within us—and around us.

But if this is the true “edification” of the believer, it is easy to understand that not very many are inclined to desire it, to pray for it. For it is costly. It costs humiliation, sacrifice, money.

Jesus says: “Take my yoke upon you and learn from me, for I am gentle and humble in heart, and you will find rest for your souls: for my yoke is beneficial and my burden is light.”

Here he says, therefore, that “his yoke” is the soul’s rest; that edification, rest for the soul, is obtained by taking “his yoke,” by learning from him to labor for his cause with something of that mind to which it is “food to do the Father’s will.”

Whether we will admit it or not, it is nevertheless true that we are all tempted to settle ourselves comfortably on earth and to trouble ourselves as little as possible about the crying need in which our brothers in the world find themselves. The “edification” which we here, as God’s people, most need when we gather around his Word at our larger or smaller meetings is this: that we receive more desire and more strength to walk fearlessly on our fixed way, in that life whose wealth and happiness are the joyful self-sacrifice for Jesus Christ’s great, world-liberating task; so that concerning us also the familiar but often forgotten word may become reality, may become truth: “As he is, so are we in this world.”

My meaning, therefore, is that the believers have received true edification during a meeting when they leave the meeting with deeper sorrow over their own powerlessness and wretchedness, with a firmer resolve to cling to grace more childlike and trustfully, and with a more burning longing to be among the ranks of those who have gone out to conquer the world for Christ. The believers are edified when, through the working of the Word, they leave a meeting fully conscious that they, in the one journey they have through life, will follow Jesus Christ.

It will, of course, cost—perhaps much: “And he said to all: If anyone wishes to come after me, he must deny himself and every day take up his cross and follow me. For whoever wishes to save his life will lose it, but whoever loses his life for my sake, he will save it.”

And it may indeed happen that the heart will at times writhe in pain over what it will cost. But nevertheless—the same heart that found salvation in Jesus says to him in full earnest:

\noindent\textit{Hymn reference:}
\href{https://hymnary.org/text/din_o_jesus_din_at_vaere}
{Hymnary.org, ``Din, O Jesus, din at være'' by Theodor Wilhelm Oldenburg.}

\smallskip

\noindent\textit{Editorial note: E.B. quotes stanzas 1 and 4 of this hymn. Because the complete hymn would have been familiar within the Norwegian Lutheran context of the period, stanzas 2, 3, and 5 are supplied here to restore the context available to the original reader.}


\begin{quote}
\medskip
\textbf{1}

Yours, O Jesus, yours to be,\\
Is my joy and all my honor,\\
Yours in life and yours in death;\\
Therefore I will struggle onward,\\
In your footsteps I will journey,\\
Finding sweet the hour of battle.

\medskip

2

Yes, I know that only he can\\
Find the crown of glory with you,\\
Glory’s crown and peace of life,\\
Who, despite all dangers of life,\\
All the foe’s deceit and darkness,\\
Faithfully walks in your steps.

\medskip

3

Underneath the cross I’ll enter,\\
Bearing it with quiet gladness,\\
That cross, Savior, which is yours;\\
Though my path leads through dark valleys,\\
Still my goal in heaven’s dwelling\\
I behold with joyful eyes.

\medskip

\textbf{4}

Lead me where death itself threatens,\\
Lead me through affliction’s fire,\\
Through the storm upon the wild sea;\\
As you will, O dearest Jesus,\\
If only you, yes, you will be\\
My own strong supporting staff.

\medskip

5

Shining in the light of your grace,\\
Lead me on the road toward the goal\\
While I wander here on earth;\\
And when all the cross is ended,\\
Let me enter into heaven,\\
Where you dwell in glorious light.

\end{quote}

And while we, according to the grace we receive, fight the fight against ourselves and our own wretchedness and against the evil in the world around us, our heart is lifted up in blessed joy at possessing God’s unchangeable love in Christ, and we count everything as loss, if only we may keep it. And he himself has taught us that when the battle has been fought to its end and the way of the cross in following him here in this foreign land has been traversed, we will then accompany him into the land of perfection and there be together with him through the long, joy-bringing “Sabbath rest that remains for God’s people.”

“Where I am, there my servant will also be.”